% Einheit 2 von 3 – Auswahlen und Binomialkoeffizient (bank/kombinatorik/e2.jsonl, zusammenbau v0.1)
%% TODO Zeitmarke Einheit 2: Marken-Zeile nicht lesbar
%% TODO Zweigzeile Teil 1: was hier gelernt wird (Satz in Schülersprache aus dem Einheitstitel) – Einheit 2
\einheitenkopf[e2]{Einheit 2 von 3 · Auswahlen und Binomialkoeffizient}
\zweigzeile{Abitur GK}
\setcounter{aufgabe}{9}

%% TODO Titel: Ich-kann-Satz fehlt in der Bank (Platzhalter: Kettenname) – Nr. 10
%% TODO Anweisung: ein Satz über der ersten Teilaufgabe fehlt in der Bank – Nr. 10
\begin{aufgabe}{Auswahlen und Binomialkoeffizient}
%% TODO \rechenplatz{n}: Schrittzahl des Grundfalls fehlt in der Bank (nur bei fester Schreibform, 2.3 b)
\begin{teile}
\teil Für einen Eisbecher werden drei verschiedene Sorten aus elf Sorten gewählt. Zählt die Reihenfolge der Sorten? \\ \kreuz{ja} \\ \kreuz{nein}
\teil Aus den vier Kindern A, B, C, D werden zwei für den Tafeldienst ausgewählt. Schreibe alle Auswahlen auf und vergleiche mit $(4 \text{ über } 2)$. \leerfeld
\teil Bei einer Ziehung werden sechs von 40 Zahlen ohne Reihenfolge gezogen. Wie viele Ziehungen gibt es? \leerfeld
\teil Vergleiche $(9 \text{ über } 2)$ und $(9 \text{ über } 7)$ ohne Rechner und gib den Wert an. \leerfeld
\teil Ein Bäcker hat neun Brotsorten. Vier kommen ins Wochenangebot, und jede davon wird an einem von vier Tagen besonders beworben. Wie viele Möglichkeiten gibt es für die Auswahl, und wie viele für die Verteilung auf die Tage \hfill (FHR 2019)? Auswahl \leerfeld , Verteilung \leerfeld
\teil Aus fünf Mädchen und vier Jungen wird ein Dreierteam mit mindestens zwei Mädchen gebildet. Wie viele Teams gibt es? \leerfeld
\teil Ein Kennwort besteht aus sechs verschiedenen Zeichen aus einem Satz von 50 Zeichen. Es enthält den Namen Lina in dieser Reihenfolge und Schreibung, dazu zwei weitere Zeichen an beliebigen Stellen, etwa L7in\#a. Wie viele solche Kennwörter gibt es \hfill (Abitur 2024 LK)? \leerfeld
\teil Sträuße aus neun Rosen werden aus den Farben Weiß, Rosa, Rot und Gelb gebunden; jeder Strauß hat genau zwei Farben. Die Anzahl der Sträuße ist $(4 \text{ über } 2) \cdot 8$. Deute beide Faktoren im Sachzusammenhang \hfill (Abitur 2020 LK).
\teil Eine sehr große Kiste enthält Teebeutel von vier Sorten. Ein Kunde greift rein zufällig sechs Beutel heraus. Wie viele Zusammenstellungen der Sorten sind möglich \hfill (FHR 2026)? \leerfeld
\end{teile}
\end{aufgabe}

%% TODO Titel: Ich-kann-Satz fehlt in der Bank (Platzhalter: Kettenname) – Nr. 11
%% TODO Anweisung: ein Satz über der ersten Teilaufgabe fehlt in der Bank – Nr. 11
\begin{aufgabe}{Auswahlen und Binomialkoeffizient – Fehler finden, Begründen, Anwendung}
\begin{teile}
\teil Aus acht Spielern werden zwei für einen Test ausgewählt. Ole rechnet so. Finde den Fehler und rechne richtig. \rechnung{8 \cdot 7 &= 56}
\teil Begründe ohne Rechnung, dass es gleich viele Möglichkeiten gibt, aus einer Gruppe $k$ Personen auszuwählen wie $n - k$ Personen.
\teil Bei einem Rubbellos sind 25 Felder, vier davon zeigen einen Gewinn. Man darf vier Felder freirubbeln. Wie groß ist die Chance, genau die vier Gewinnfelder zu treffen? 1 zu \leerfeld
\end{teile}
\end{aufgabe}
